\section{Conclusion}
\label{sec:conclusion}

Forward fine-tuning can create a hidden failure mode in code LLMs. In our
controlled testbed, the adaptation step that raises forward obfuscation accuracy from below 5\%
to between 12.0\% and 81.3\% before self-correction,  depending on transformation type, drives down reverse deobfuscation accuracy to 0\%. The reverse
outputs remain executable and semantically valid, yet they stay near-verbatim
reproductions of the obfuscated input rather than steps back toward the original
readable program. 
% Forward success therefore does not entail preservation of the
% inverse capability.
Forward success, therefore, does not preserve the inverse capability.

The recovery experiments qualify that claim. Balanced bidirectional supervision
restores much of the reverse performance on the lexical and structural tasks,
while contrastive fine-tuning restores the lexical inverse without reverse
labels. Literal Encryption resists all three strategies. The translation and
compilation settings reproduce the directional trend outside the primary
obfuscation study. When a fine-tuned code LLM is destined for a setting in which
the inverse transformation matters, that inverse should be evaluated in its own
right, and BRF exposes the capability loss that forward-only benchmarks miss.

\noindent\textbf{Data Availability.} Code, datasets, and evaluation results are available at
\url{https://anonymous.4open.science/r/Bidirectional_Reasoning_Framework-B111/README.md}
